\section{指标定义：原始、条件与联合量}\label{app:metrics}
理想目标纯态$\ket\psi$的原始态重叠是$F=\bra\psi\rho\ket\psi$。局域down比较CLR约化态，up及往返比较CI；合法过程块保留完整四线目标像。若四维块K按目标/入口态排序，则
\begin{equation}
 \Fe^{\rm legal}=\frac{|\tr K|^2}{16},\qquad
 P_{\rm image}=\frac14\tr(K^\dagger K).
\end{equation}
前者也检查相位，后者只看目标像人口，均未按列归一。

完整query的全输出Choi态是$J_{\rm full}=(I\otimes\mathcal E)(\proj{\Phi_8})$；输出限制AD的01块后为J，迹$S\leq1$，泄露时它不归一。Haar平均原始目标态重叠为
\begin{equation}
 \overline F_{\rm Haar}=\frac{8\Fe+S}{9}.
\end{equation}
只有S为1才退化成常用迹保持形式。

给定输入，真实辅助零投影$P_{\mathcal R0}$定义
\begin{equation}
 p_0=\tr(P_{\mathcal R0}\rho),\qquad
 \rho_0=P_{\mathcal R0}\rho P_{\mathcal R0}/p_0.
\end{equation}
目标地址测量为$\Pi_a$时，三种成功指标分别为
\begin{align}
 p_{\rm raw}&=\tr(\Pi_a\rho),\\
 p_{\rm cond}&=\tr(\Pi_a\rho_0),\\
 p_{\rm joint}&=\tr(\Pi_aP_{\mathcal R0}\rho P_{\mathcal R0}).
\end{align}
Choi平均接受率与某算法输入接受率不同，不能混用。

\clearpage\section{冻结参考参数与原生单门表}\label{app:parameters}
以下来自已有模拟JSON，计算用全精度，表中四舍五入。参考标签q111/q121/q130/q112中q121是中心I；不是完整17角色硬件映射，也不是当前硬件上传值。

附录末尾的四份局域门表中，$\phi_{\rm nom}$为协议相位；$\phi_{\rm buf}$为本地帧从零开始、累计补偿后的驱动轴，折回$[-\pi,\pi]$。它不含绝对载波，也不能直接当全局AWG相位。F类更新时长0；同层异位单门并行，层间隔离。
\subsection{R5的配对矩阵与实现}
\begin{equation}V_{\mathrm{down}}=\begin{pmatrix}0.923048+0.000000i & -0.001145-0.384684i\\0.274545-0.269459i & 0.648521+0.656839i\end{pmatrix}.\end{equation}
\begin{equation}V_{\mathrm{up}}=\begin{pmatrix}0.922191+0.000000i & 0.273149-0.273776i\\0.000251-0.386734i & 0.652412+0.651763i\end{pmatrix}.\end{equation}

对$V=e^{i\alpha}R_z(\beta)R_y(\theta)R_z(\gamma)$定义
\begin{equation}
 \widetilde Z(\delta):F_{01}(-\delta)\seq F_{12}(\delta/2).
\end{equation}
在012空间它等于嵌入$R_z(\delta)$乘公共相位。V的时间展开是
\begin{equation}
\begin{aligned}
 &\widetilde Z(\gamma)\seq\widetilde Z(\pi/2)\seq Y_{01}/2\seq\widetilde Z(\theta)
 \seq(-Y_{01}/2)\\
 &\seq\widetilde Z(-\pi/2)\seq\widetilde Z(\beta)\seq F_{12}(\alpha).
\end{aligned}
\end{equation}
末尾$F_{12}(\alpha)$维持01块相对2态的相位，不能丢弃。$V^\dagger$按矩阵重新分解；更高能级作用由原生展开与脉冲决定，不另假设恒等。

\begin{table}[htbp]\centering\small
\caption{单门冻结参数（程序幅度单位换算为MHz）}\label{tab:singleparams}
\begin{tabular}{@{}llrrrr@{}}\toprule
线 & 门 & T/ns & A/MHz & $\delta$/MHz & DRAG\\\midrule
C & $X_{01}/2$ & 32 & 15.623205 & -0.190561 & 0.144406\\
C & $X_{01}$ & 32 & 31.195108 & -1.076721 & 0.134111\\
C & $X_{12}$ & 40 & 17.692091 & -0.538408 & 0.082153\\
I & $X_{01}/2$ & 32 & 15.623236 & -0.190544 & 0.144409\\
I & $X_{01}$ & 32 & 31.156308 & -1.373690 & 0.149271\\
I & $X_{12}$ & 40 & 17.692086 & -0.538369 & 0.082149\\
L & $X_{01}/2$ & 32 & 15.623216 & -0.190535 & 0.144403\\
L & $X_{01}$ & 32 & 31.257109 & -0.507768 & 0.105243\\
L & $X_{12}$ & 40 & 17.692085 & -0.538399 & 0.082151\\
R & $X_{01}/2$ & 32 & 15.623367 & -0.180091 & 0.145040\\
R & $X_{01}$ & 32 & 31.271875 & -0.296084 & 0.095514\\
R & $X_{12}$ & 40 & 17.688586 & -0.545110 & 0.085622\\
\bottomrule\end{tabular}\end{table}
\begin{table}[htbp]\centering\small
\caption{CZ冻结参数，三边峰值因子均为1}\label{tab:czparams}
\begin{tabular}{@{}llrrr@{}}\toprule
边 & 调谐线 & $f^{\CZ}$/GHz & T/ns & 坡/ns\\\midrule
C--I & I & 4.818937693 & 30.641834 & 5.725930\\
I--L & L & 5.244278591 & 56.371311 & 11.894511\\
I--R & R & 5.259764008 & 64.820457 & 16.500000\\
\bottomrule\end{tabular}\end{table}


\begin{table}[htbp]\centering\small
\caption{端点帧：down放输出，up放输入，角度单位rad}\label{tab:endpoint}
\begin{tabular}{@{}lrrr@{}}\toprule
方案/方向 & $\phi_C$ & $\phi_L$ & $\phi_R$\\\midrule
R5 down & 2.136815 & -0.134507 & 2.348115\\
R5 up & 2.159159 & -0.153203 & 2.334839\\
R6 down & 2.139030 & -0.133450 & 2.956958\\
R6 up & 2.114272 & -0.143818 & 3.006746\\
\bottomrule\end{tabular}\end{table}


\clearpage{\small\setlength{\tabcolsep}{4pt}\fontsize{9}{10.6}\selectfont\renewcommand{\arraystretch}{1.0}
\begin{longtable}{@{}rrrllrrr@{}}
\caption{R5 down：18层，713.025368 ns}\label{tab:r5down}\\\toprule
序 & 层 & 起点/ns & 门 & 线 & $\phi_{\rm nom}$ & $\phi_{\rm buf}$ & 门长/ns\\\midrule\endfirsthead
\multicolumn{8}{c}{续表：R5 down：18层，713.025368 ns}\\\toprule
序 & 层 & 起点/ns & 门 & 线 & $\phi_{\rm nom}$ & $\phi_{\rm buf}$ & 门长/ns\\\midrule\endhead
\midrule\multicolumn{8}{r}{下页续}\endfoot
\bottomrule\endlastfoot
1 & 1 & 0.000 & $X_{01}$ & I & 0.00000 & 0.00000 & 32.000\\
2 & 2 & 32.000 & $X_{12}$ & I & 0.00000 & -0.01213 & 40.000\\
3 & 3 & 72.000 & $X_{01}/2$ & I & -1.57080 & -1.46338 & 32.000\\
4 & 4 & 104.000 & $\CZ$ & C,I & --- & --- & 30.642\\
5 & 5 & 134.642 & $X_{01}/2$ & I & 1.57080 & -1.49542 & 32.000\\
6 & 6 & 166.642 & $X_{12}$ & I & 0.00000 & -3.06430 & 40.000\\
7 & 7 & 206.642 & $X_{01}$ & I & 0.00000 & 3.08650 & 32.000\\
8 & 7 & 206.642 & $F_{01}$ & I & -1.56782 & --- & 0.000\\
9 & 7 & 206.642 & $F_{12}$ & I & 0.78391 & --- & 0.000\\
10 & 7 & 206.642 & $F_{01}$ & I & -1.57080 & --- & 0.000\\
11 & 7 & 206.642 & $F_{12}$ & I & 0.78540 & --- & 0.000\\
12 & 8 & 238.642 & $X_{01}/2$ & I & 1.57080 & 1.79488 & 32.000\\
13 & 8 & 238.642 & $F_{01}$ & I & -0.78973 & --- & 0.000\\
14 & 8 & 238.642 & $F_{12}$ & I & 0.39487 & --- & 0.000\\
15 & 9 & 270.642 & $X_{01}/2$ & I & -1.57080 & -2.09813 & 32.000\\
16 & 9 & 270.642 & $F_{01}$ & I & 1.57080 & --- & 0.000\\
17 & 9 & 270.642 & $F_{12}$ & I & -0.78540 & --- & 0.000\\
18 & 9 & 270.642 & $F_{01}$ & I & 0.77605 & --- & 0.000\\
19 & 9 & 270.642 & $F_{12}$ & I & -0.38802 & --- & 0.000\\
20 & 9 & 270.642 & $F_{12}$ & I & 0.39589 & --- & 0.000\\
21 & 9 & 270.642 & $X_{01}/2$ & L & -1.57080 & -1.57080 & 32.000\\
22 & 10 & 302.642 & $\CZ$ & I,L & --- & --- & 56.371\\
23 & 11 & 359.013 & $X_{01}/2$ & L & 1.57080 & -2.58760 & 32.000\\
24 & 11 & 359.013 & $X_{01}/2$ & I & -1.57080 & 0.54838 & 32.000\\
25 & 12 & 391.013 & $\CZ$ & L,I & --- & --- & 56.371\\
26 & 13 & 447.384 & $X_{01}/2$ & I & 1.57080 & -2.29354 & 32.000\\
27 & 13 & 447.384 & $F_{01}$ & L & 2.36554 & --- & 0.000\\
28 & 13 & 447.384 & $F_{12}$ & L & -1.18277 & --- & 0.000\\
29 & 13 & 447.384 & $F_{01}$ & L & -1.57080 & --- & 0.000\\
30 & 13 & 447.384 & $F_{12}$ & L & 0.78540 & --- & 0.000\\
31 & 13 & 447.384 & $X_{01}/2$ & L & 1.57080 & 0.33193 & 32.000\\
32 & 13 & 447.384 & $F_{01}$ & L & -0.78973 & --- & 0.000\\
33 & 13 & 447.384 & $F_{12}$ & L & 0.39487 & --- & 0.000\\
34 & 14 & 479.384 & $X_{01}/2$ & L & -1.57080 & 2.72210 & 32.000\\
35 & 14 & 479.384 & $F_{01}$ & L & 1.57080 & --- & 0.000\\
36 & 14 & 479.384 & $F_{12}$ & L & -0.78540 & --- & 0.000\\
37 & 14 & 479.384 & $F_{01}$ & L & -1.57377 & --- & 0.000\\
38 & 14 & 479.384 & $F_{12}$ & L & 0.78689 & --- & 0.000\\
39 & 14 & 479.384 & $F_{12}$ & L & -0.39589 & --- & 0.000\\
40 & 14 & 479.384 & $X_{12}$ & I & 0.00000 & -1.18282 & 40.000\\
41 & 14 & 479.384 & $X_{01}/2$ & R & -1.57080 & -1.57080 & 32.000\\
42 & 15 & 519.384 & $\CZ$ & I,R & --- & --- & 64.820\\
43 & 16 & 584.205 & $X_{01}/2$ & R & 1.57080 & 0.44708 & 32.000\\
44 & 16 & 584.205 & $X_{01}/2$ & I & -1.57080 & 0.97514 & 32.000\\
45 & 17 & 616.205 & $\CZ$ & R,I & --- & --- & 64.820\\
46 & 18 & 681.025 & $X_{01}/2$ & I & 1.57080 & -1.87059 & 32.000\\
47 & 18 & 681.025 & $F$ & C & 2.13681 & --- & 0.000\\
48 & 18 & 681.025 & $F$ & L & -0.13451 & --- & 0.000\\
49 & 18 & 681.025 & $F$ & R & 2.34811 & --- & 0.000\\
\end{longtable}}

\clearpage{\small\setlength{\tabcolsep}{4pt}\fontsize{9}{10.6}\selectfont\renewcommand{\arraystretch}{1.0}
\begin{longtable}{@{}rrrllrrr@{}}
\caption{R5 up：18层，713.025368 ns}\label{tab:r5up}\\\toprule
序 & 层 & 起点/ns & 门 & 线 & $\phi_{\rm nom}$ & $\phi_{\rm buf}$ & 门长/ns\\\midrule\endfirsthead
\multicolumn{8}{c}{续表：R5 up：18层，713.025368 ns}\\\toprule
序 & 层 & 起点/ns & 门 & 线 & $\phi_{\rm nom}$ & $\phi_{\rm buf}$ & 门长/ns\\\midrule\endhead
\midrule\multicolumn{8}{r}{下页续}\endfoot
\bottomrule\endlastfoot
1 & 1 & 0.000 & $F$ & C & 2.15916 & --- & 0.000\\
2 & 1 & 0.000 & $F$ & L & -0.15320 & --- & 0.000\\
3 & 1 & 0.000 & $F$ & R & 2.33484 & --- & 0.000\\
4 & 1 & 0.000 & $X_{01}/2$ & I & -1.57080 & -1.57080 & 32.000\\
5 & 2 & 32.000 & $\CZ$ & R,I & --- & --- & 64.820\\
6 & 3 & 96.820 & $X_{01}/2$ & I & 1.57080 & 1.86666 & 32.000\\
7 & 3 & 96.820 & $X_{01}/2$ & R & -1.57080 & -0.39589 & 32.000\\
8 & 4 & 128.820 & $\CZ$ & I,R & --- & --- & 64.820\\
9 & 5 & 193.641 & $X_{01}/2$ & R & 1.57080 & 1.62199 & 32.000\\
10 & 5 & 193.641 & $X_{12}$ & I & 0.00000 & 0.53957 & 40.000\\
11 & 5 & 193.641 & $F_{01}$ & L & 1.57144 & --- & 0.000\\
12 & 5 & 193.641 & $F_{12}$ & L & -0.78572 & --- & 0.000\\
13 & 5 & 193.641 & $F_{01}$ & L & -1.57080 & --- & 0.000\\
14 & 5 & 193.641 & $F_{12}$ & L & 0.78540 & --- & 0.000\\
15 & 5 & 193.641 & $X_{01}/2$ & L & 1.57080 & 1.41824 & 32.000\\
16 & 5 & 193.641 & $F_{01}$ & L & -0.79418 & --- & 0.000\\
17 & 5 & 193.641 & $F_{12}$ & L & 0.39709 & --- & 0.000\\
18 & 6 & 233.641 & $X_{01}/2$ & L & -1.57080 & -2.47922 & 32.000\\
19 & 6 & 233.641 & $F_{01}$ & L & 1.57080 & --- & 0.000\\
20 & 6 & 233.641 & $F_{12}$ & L & -0.78540 & --- & 0.000\\
21 & 6 & 233.641 & $F_{01}$ & L & -0.78654 & --- & 0.000\\
22 & 6 & 233.641 & $F_{12}$ & L & 0.39327 & --- & 0.000\\
23 & 6 & 233.641 & $F_{12}$ & L & -0.39245 & --- & 0.000\\
24 & 6 & 233.641 & $X_{01}/2$ & I & -1.57080 & -1.14785 & 32.000\\
25 & 7 & 265.641 & $\CZ$ & L,I & --- & --- & 56.371\\
26 & 8 & 322.012 & $X_{01}/2$ & I & 1.57080 & 2.29341 & 32.000\\
27 & 8 & 322.012 & $X_{01}/2$ & L & -1.57080 & 0.42982 & 32.000\\
28 & 9 & 354.012 & $\CZ$ & I,L & --- & --- & 56.371\\
29 & 10 & 410.384 & $X_{01}/2$ & L & 1.57080 & -0.58699 & 32.000\\
30 & 10 & 410.384 & $F_{01}$ & I & -2.35505 & --- & 0.000\\
31 & 10 & 410.384 & $F_{12}$ & I & 1.17752 & --- & 0.000\\
32 & 10 & 410.384 & $F_{01}$ & I & -1.57080 & --- & 0.000\\
33 & 10 & 410.384 & $F_{12}$ & I & 0.78540 & --- & 0.000\\
34 & 10 & 410.384 & $X_{01}/2$ & I & 1.57080 & -1.33276 & 32.000\\
35 & 10 & 410.384 & $F_{01}$ & I & -0.79418 & --- & 0.000\\
36 & 10 & 410.384 & $F_{12}$ & I & 0.39709 & --- & 0.000\\
37 & 11 & 442.384 & $X_{01}/2$ & I & -1.57080 & 1.05297 & 32.000\\
38 & 11 & 442.384 & $F_{01}$ & I & 1.57080 & --- & 0.000\\
39 & 11 & 442.384 & $F_{12}$ & I & -0.78540 & --- & 0.000\\
40 & 11 & 442.384 & $F_{01}$ & I & 1.57015 & --- & 0.000\\
41 & 11 & 442.384 & $F_{12}$ & I & -0.78507 & --- & 0.000\\
42 & 11 & 442.384 & $F_{12}$ & I & 0.39245 & --- & 0.000\\
43 & 12 & 474.384 & $X_{01}$ & I & 0.00000 & -0.48016 & 32.000\\
44 & 13 & 506.384 & $X_{12}$ & I & 0.00000 & 2.41640 & 40.000\\
45 & 14 & 546.384 & $X_{01}/2$ & I & -1.57080 & -1.94354 & 32.000\\
46 & 15 & 578.384 & $\CZ$ & C,I & --- & --- & 30.642\\
47 & 16 & 609.025 & $X_{01}/2$ & I & 1.57080 & -1.97559 & 32.000\\
48 & 17 & 641.025 & $X_{12}$ & I & 0.00000 & -0.63577 & 40.000\\
49 & 18 & 681.025 & $X_{01}$ & I & 0.00000 & 2.60634 & 32.000\\
\end{longtable}}

\clearpage{\small\setlength{\tabcolsep}{4pt}\fontsize{9}{10.6}\selectfont\renewcommand{\arraystretch}{1.0}
\begin{longtable}{@{}rrrllrrr@{}}
\caption{R6 down：19层，769.396679 ns}\label{tab:r6down}\\\toprule
序 & 层 & 起点/ns & 门 & 线 & $\phi_{\rm nom}$ & $\phi_{\rm buf}$ & 门长/ns\\\midrule\endfirsthead
\multicolumn{8}{c}{续表：R6 down：19层，769.396679 ns}\\\toprule
序 & 层 & 起点/ns & 门 & 线 & $\phi_{\rm nom}$ & $\phi_{\rm buf}$ & 门长/ns\\\midrule\endhead
\midrule\multicolumn{8}{r}{下页续}\endfoot
\bottomrule\endlastfoot
1 & 1 & 0.000 & $X_{01}/2$ & L & -1.57080 & -1.57080 & 32.000\\
2 & 2 & 32.000 & $\CZ$ & I,L & --- & --- & 56.371\\
3 & 3 & 88.371 & $X_{01}/2$ & L & 1.57080 & -2.58760 & 32.000\\
4 & 3 & 88.371 & $X_{01}$ & I & 0.00000 & 0.26136 & 32.000\\
5 & 4 & 120.371 & $X_{12}$ & I & 0.00000 & 0.24924 & 40.000\\
6 & 5 & 160.371 & $X_{01}/2$ & I & -1.57080 & -1.20202 & 32.000\\
7 & 6 & 192.371 & $\CZ$ & C,I & --- & --- & 30.642\\
8 & 7 & 223.013 & $X_{01}/2$ & I & 1.57080 & -1.23406 & 32.000\\
9 & 8 & 255.013 & $X_{12}$ & I & 0.00000 & -2.80294 & 40.000\\
10 & 9 & 295.013 & $X_{01}$ & I & 0.00000 & -2.93532 & 32.000\\
11 & 10 & 327.013 & $X_{01}/2$ & I & -1.57080 & 2.05326 & 32.000\\
12 & 11 & 359.013 & $\CZ$ & L,I & --- & --- & 56.371\\
13 & 12 & 415.384 & $X_{01}/2$ & I & 1.57080 & -0.78866 & 32.000\\
14 & 13 & 447.384 & $X_{12}$ & I & 0.00000 & -2.39392 & 40.000\\
15 & 13 & 447.384 & $X_{01}/2$ & L & -1.57080 & 2.67878 & 32.000\\
16 & 14 & 487.384 & $\CZ$ & I,L & --- & --- & 56.371\\
17 & 15 & 543.756 & $X_{01}/2$ & L & 1.57080 & 1.66197 & 32.000\\
18 & 15 & 543.756 & $X_{01}/2$ & R & -1.57080 & -1.57080 & 32.000\\
19 & 16 & 575.756 & $\CZ$ & I,R & --- & --- & 64.820\\
20 & 17 & 640.576 & $X_{01}/2$ & R & 1.57080 & 0.44708 & 32.000\\
21 & 17 & 640.576 & $X_{01}/2$ & I & -1.57080 & 2.74138 & 32.000\\
22 & 18 & 672.576 & $\CZ$ & R,I & --- & --- & 64.820\\
23 & 19 & 737.397 & $X_{01}/2$ & I & 1.57080 & -0.10435 & 32.000\\
24 & 19 & 737.397 & $F$ & C & 2.13903 & --- & 0.000\\
25 & 19 & 737.397 & $F$ & L & -0.13345 & --- & 0.000\\
26 & 19 & 737.397 & $F$ & R & 2.95696 & --- & 0.000\\
\end{longtable}}

\clearpage{\small\setlength{\tabcolsep}{4pt}\fontsize{9}{10.6}\selectfont\renewcommand{\arraystretch}{1.0}
\begin{longtable}{@{}rrrllrrr@{}}
\caption{R6 up：19层，769.396679 ns}\label{tab:r6up}\\\toprule
序 & 层 & 起点/ns & 门 & 线 & $\phi_{\rm nom}$ & $\phi_{\rm buf}$ & 门长/ns\\\midrule\endfirsthead
\multicolumn{8}{c}{续表：R6 up：19层，769.396679 ns}\\\toprule
序 & 层 & 起点/ns & 门 & 线 & $\phi_{\rm nom}$ & $\phi_{\rm buf}$ & 门长/ns\\\midrule\endhead
\midrule\multicolumn{8}{r}{下页续}\endfoot
\bottomrule\endlastfoot
1 & 1 & 0.000 & $F$ & C & 2.11427 & --- & 0.000\\
2 & 1 & 0.000 & $F$ & L & -0.14382 & --- & 0.000\\
3 & 1 & 0.000 & $F$ & R & 3.00675 & --- & 0.000\\
4 & 1 & 0.000 & $X_{01}/2$ & I & -1.57080 & -1.57080 & 32.000\\
5 & 2 & 32.000 & $\CZ$ & R,I & --- & --- & 64.820\\
6 & 3 & 96.820 & $X_{01}/2$ & I & 1.57080 & 1.86666 & 32.000\\
7 & 3 & 96.820 & $X_{01}/2$ & R & -1.57080 & 0.27602 & 32.000\\
8 & 4 & 128.820 & $\CZ$ & I,R & --- & --- & 64.820\\
9 & 5 & 193.641 & $X_{01}/2$ & R & 1.57080 & 2.29389 & 32.000\\
10 & 5 & 193.641 & $X_{01}/2$ & L & -1.57080 & -1.71461 & 32.000\\
11 & 6 & 225.641 & $\CZ$ & I,L & --- & --- & 56.371\\
12 & 7 & 282.012 & $X_{01}/2$ & L & 1.57080 & -2.73142 & 32.000\\
13 & 7 & 282.012 & $X_{12}$ & I & 0.00000 & 0.80093 & 40.000\\
14 & 8 & 322.012 & $X_{01}/2$ & I & -1.57080 & -0.88649 & 32.000\\
15 & 9 & 354.012 & $\CZ$ & L,I & --- & --- & 56.371\\
16 & 10 & 410.384 & $X_{01}/2$ & I & 1.57080 & 2.55477 & 32.000\\
17 & 11 & 442.384 & $X_{01}$ & I & 0.00000 & 1.02229 & 32.000\\
18 & 12 & 474.384 & $X_{12}$ & I & 0.00000 & 1.20994 & 40.000\\
19 & 13 & 514.384 & $X_{01}/2$ & I & -1.57080 & -0.44109 & 32.000\\
20 & 14 & 546.384 & $\CZ$ & C,I & --- & --- & 30.642\\
21 & 15 & 577.025 & $X_{01}/2$ & I & 1.57080 & -0.47313 & 32.000\\
22 & 16 & 609.025 & $X_{12}$ & I & 0.00000 & -1.84223 & 40.000\\
23 & 17 & 649.025 & $X_{01}$ & I & 0.00000 & -2.17439 & 32.000\\
24 & 17 & 649.025 & $X_{01}/2$ & L & -1.57080 & 2.53496 & 32.000\\
25 & 18 & 681.025 & $\CZ$ & I,L & --- & --- & 56.371\\
26 & 19 & 737.397 & $X_{01}/2$ & L & 1.57080 & 1.51815 & 32.000\\
\end{longtable}}


\section{工程组织与复现}\label{app:reproduce}
主文件\code{main.tex}，正文分为\code{sections/}内五个源文件；\code{images/}含25张矢量图，\code{tables/}含JSON导出的数表，\code{data/}保留主要结果与冻结参数。XeLaTeX两遍可解析目录与引用，无需图片公式。

\code{supplement/}保存R5/R6完整整网逐层Markdown及调度JSON，另有局域名义与缓冲相位JSON。主文宏观时间表和局域门表用于讲清逻辑，数百行完整执行表单独提供以便核对。

上一级\code{prepare_latex_report.py}只导出数表、复制图及记录指纹，不下发实验、不重跑传播。正文源文件不被该脚本覆盖。完整数值入口、依赖和原始报告见上一级\code{README.md}。本版保留旧PDF，单独提供原生公式和重构叙述。
